在定理 \ref{thm:conclusion-godel-semidirect-law} 与命题 \ref{prop:conclusion-prime-register-elliptic-representation} 的接口之外，还存在一类把“伸缩轴”与“平移外置”统一封装为二维线性对象的仿射正规形。
该正规形的关键特征是：伸缩在主对角上以乘法聚合，而平移以半直积律携带次序敏感的进位结构；其几何像为单位圆的仿射椭圆族，参数可由横截距与切线斜率直接反演。

\begin{theorem}[伸缩--平移外置的仿射正规形与半直积律]\label{thm:conclusion-affine-normal-form-semidirect}
对 \(N\in\NN_{>0}\)、\(k\in\NN\) 定义
\[
A_{N,k}:=
\begin{pmatrix}
\sqrt N & \dfrac{k}{\sqrt N}\\[2mm]
0 & \dfrac{1}{\sqrt N}
\end{pmatrix}\in \mathrm{SL}_2(\RR),
\qquad
R:=A_{1,1}=
\begin{pmatrix}
1&1\\
0&1
\end{pmatrix},
\qquad
D_p:=A_{p,0}=
\begin{pmatrix}
\sqrt p&0\\
0&1/\sqrt p
\end{pmatrix}\ (p\ \text{素数}).
\]
令 \(\mathcal A^{+}\) 为由 \(\{R\}\cup\{D_p:\ p\ \text{素数}\}\) 在 \(\mathrm{SL}_2(\RR)\) 中生成的子幺半群（仅允许非负幂与有限乘积）。
则存在显式同构
\[
\Phi:\NN_{>0}\ltimes \NN\ \xrightarrow{\ \cong\ }\ \mathcal A^{+},
\qquad
\Phi(N,k)=A_{N,k},
\]
其中半直积乘法为
\[
(N,k)\cdot(M,\ell)=(NM,\ k+N\ell).
\]
等价地，
\[
A_{N,k}A_{M,\ell}=A_{NM,\ k+N\ell}\qquad(\forall\,N,M\in\NN_{>0},\ k,\ell\in\NN),
\]
并且 \(\mathcal A^{+}\) 的每个元素在坐标 \((N,k)\) 下表出唯一。
\end{theorem}
\begin{proof}
直接计算得
\[
\begin{pmatrix}\sqrt N & k/\sqrt N\\ 0&1/\sqrt N\end{pmatrix}
\begin{pmatrix}\sqrt M & \ell/\sqrt M\\ 0&1/\sqrt M\end{pmatrix}
=
\begin{pmatrix}\sqrt{NM} & \dfrac{k+N\ell}{\sqrt{NM}}\\[2mm] 0 & \dfrac1{\sqrt{NM}}\end{pmatrix}
=A_{NM,k+N\ell},
\]
从而 \(\{A_{N,k}\}\) 对乘法封闭并包含生成元 \(R=A_{1,1}\) 与 \(D_p=A_{p,0}\)，故 \(\mathcal A^{+}\subseteq \{A_{N,k}\}\)。
反向包含由归纳与分解 \(A_{N,k}=R^{k}D_N\)、\(D_N=\prod_{p}D_p^{v_p(N)}\) 立即成立。
唯一性由矩阵条目给出：若 \(A_{N,k}=A_{N',k'}\)，则左上角推出 \(N=N'\)，继而右上角推出 \(k=k'\)。
\leanverified{A\_N\_k}
\leanverified{semidirectMul}
\leanverified{sqrt\_pos\_of\_pos}
\leanverified{sqrt\_mul\_sqrt}
\leanverified{A\_N\_k\_mul}
\leanverified{semidirectMul\_def}
\leanverified{paper\_conclusion\_affine\_normal\_form\_semidirect}
\end{proof}

\begin{theorem}[primorial 分层进位与素数伸缩门的可逆外置]\label{thm:conclusion-primorial-mixed-radix-affine}
令 \(p_1<p_2<\cdots\) 为素数序列并定义 primorial 权重
\[
P_0:=1,\qquad P_t:=\prod_{j=1}^{t}p_j\quad(t\ge 1).
\]
对任意长度 \(T\) 的分层素数位数码 \(a=(a_1,\dots,a_T)\)（满足 \(0\le a_t<p_t\)），定义
\[
K(a):=\sum_{t=1}^{T}a_t\,P_{t-1}\in\NN.
\]
则：
\begin{enumerate}
  \item 映射 \(a\mapsto K(a)\) 给出集合双射
  \[
  \{(a_1,\dots,a_T):0\le a_t<p_t\}\ \xrightarrow{\ \cong\ }\ \{0,1,\dots,P_T-1\}.
  \]
  \item 定义
  \[
  G(a):=R^{a_1}D_{p_1}\,R^{a_2}D_{p_2}\cdots R^{a_T}D_{p_T}\in \mathrm{SL}_2(\RR),
  \]
  则有显式恒等式
\[
  G(a)=A_{P_T,\ K(a)}.
  \]
\end{enumerate}
\leanverified{paper\_conclusion\_primorial\_mixed\_radix\_affine}
\leanverified{paper\_conclusion\_primorial\_mixed\_radix\_affine\_seeds}
\end{theorem}
\begin{proof}
双射性为混合进位表示的标准分解：对任意 \(0\le k<P_T\)，递归定义
\[
k_1:=k,\qquad a_t:=k_t\bmod p_t,\qquad k_{t+1}:=\frac{k_t-a_t}{p_t}\qquad(1\le t\le T),
\]
则 \(0\le a_t<p_t\)，且由 \(k<P_T\) 得 \(k_{T+1}=0\)，反向代入得到 \(k=K(a)\)，从而 \(K\) 为双射。
\par
矩阵恒等式由定理 \ref{thm:conclusion-affine-normal-form-semidirect} 的乘法律与 \(R^{m}=A_{1,m}\) 归纳得到：
\[
G(a)
=A_{p_1,a_1}A_{p_2,a_2}\cdots A_{p_T,a_T}
=A_{P_T,\ a_1+p_1a_2+\cdots+P_{T-1}a_T}
=A_{P_T,K(a)}.
\]
\end{proof}

\begin{corollary}[乘法 Gödel 编码与加法 primorial 编码的互逆]\label{cor:conclusion-primorial-additive-vs-godel-multiplicative}
在定理 \ref{thm:conclusion-primorial-mixed-radix-affine} 的设定下，定义乘法式 Gödel 数
\[
G_{\times}(a):=\prod_{t=1}^{T}p_t^{a_t+1}\in\NN_{>0}.
\]
则存在显式互逆链
\[
a\ \longleftrightarrow\ K(a)\ \longleftrightarrow\ G_{\times}(a),
\]
其中 \(K(a)\to a\) 由逐层取模与整除给出，而 \(G_{\times}(a)\to a\) 由素因子分解得到 \(a_t=v_{p_t}(G_{\times}(a))-1\)。
\end{corollary}
\begin{proof}
由算术基本定理，\(\{v_{p_t}(G_{\times}(a))\}\) 唯一确定 \(a\)；而 \(K(a)\) 的混合进位解码递归同样唯一恢复 \(a\)。
两者均与 \(a\) 双射，从而互可恢复。
\leanverified{paper\_conclusion\_primorial\_additive\_vs\_godel\_multiplicative\_small}
\leanverified{paper\_conclusion\_primorial\_additive\_vs\_godel\_multiplicative}
\leanverified{mixedRadixVal}
\leanverified{godelMul}
\leanverified{godelMul\_injective\_zero}
\leanverified{godelMul\_injective\_one}
\leanverified{mixedRadixVal\_injective\_zero}
\leanverified{mixedRadixVal\_injective\_one}
\leanverified{godelMul\_exponent\_eq}
\leanverified{godelMul\_injective\_coprime}
\end{proof}

\begin{proposition}[Gödel--仿射椭圆的代数方程与参数可识别性]\label{prop:conclusion-godel-affine-ellipse-equation-identifiability}
\leanverified{paper\_conclusion\_godel\_affine\_ellipse\_equation\_identifiability}
对任意 \(N\in\NN_{>0}\)、\(k\in\NN\)，定义仿射椭圆
\[
\mathcal E_{N,k}:=A_{N,k}(S^1)\subset \RR^2.
\]
则：
\begin{enumerate}
  \item \(\mathcal E_{N,k}\) 的显式代数方程为
  \[
  \frac{(x-ky)^2}{N}+Ny^2=1.
  \]
  \item 椭圆 \(\mathcal E_{N,k}\) 唯一确定参数 \((N,k)\)。
  更精确地，横截距满足
  \[
  N=\Bigl(\max\{|x|:(x,0)\in\mathcal E_{N,k}\}\Bigr)^2,
  \]
  且当 \(k\neq 0\) 时，点 \((\sqrt N,0)\) 处的切线斜率为 \(1/k\)；当 \(k=0\) 时该点切线为竖直线 \(x=\sqrt N\)。
\end{enumerate}
\end{proposition}
\begin{proof}
令 \(v=(x,y)^{\mathsf T}\)。由 \(\mathcal E_{N,k}=A_{N,k}(S^1)\) 得
\(
v\in\mathcal E_{N,k}\iff \|A_{N,k}^{-1}v\|^2=1
\)，即
\[
v^{\mathsf T}\bigl(A_{N,k}^{-\mathsf T}A_{N,k}^{-1}\bigr)v=1.
\]
计算
\[
A_{N,k}^{-1}=
\begin{pmatrix}
1/\sqrt N&-k/\sqrt N\\
0&\sqrt N
\end{pmatrix},
\qquad
A_{N,k}^{-\mathsf T}A_{N,k}^{-1}
=
\begin{pmatrix}
1/N & -k/N\\
-k/N & k^2/N+N
\end{pmatrix},
\]
展开即得所示方程。
\par
取 \(y=0\) 得 \(x=\pm\sqrt N\)，故 \(N\) 由横截距唯一恢复。
对隐函数 \(F(x,y)=\frac{(x-ky)^2}{N}+Ny^2-1\) 有 \(\frac{\dd y}{\dd x}=-F_x/F_y\)。
在 \((x,y)=(\sqrt N,0)\) 处，
\(
F_x=2/\sqrt N
\) 且
\(
F_y=-2k/\sqrt N
\)，故当 \(k\neq 0\) 有 \(\dd y/\dd x=1/k\)；当 \(k=0\) 时 \(F_y=0\) 给出竖直切线。
\end{proof}

\leanverified{paper\_conclusion\_primorial\_ellipse\_unique\_factorization}
\begin{theorem}[分层外置门的唯一分解与椭圆族充满性]\label{thm:conclusion-primorial-ellipse-unique-factorization}
固定长度 \(T\) 与素数序列 \((p_1,\dots,p_T)\)。
定义基本外置门
\[
g_t(a):=R^{a}D_{p_t}\qquad (0\le a<p_t),
\]
并定义
\[
g(a_1,\dots,a_T):=g_1(a_1)\,g_2(a_2)\cdots g_T(a_T),
\qquad
\mathcal E(a_1,\dots,a_T):=g(a_1,\dots,a_T)(S^1).
\]
则映射 \((a_1,\dots,a_T)\mapsto \mathcal E(a_1,\dots,a_T)\) 为单射，且其像恰为集合
\[
\{\mathcal E_{P_T,k}:\ 0\le k<P_T\}.
\]
\end{theorem}
\begin{proof}
由定理 \ref{thm:conclusion-primorial-mixed-radix-affine} 得
\(
g(a_1,\dots,a_T)=A_{P_T,K(a)}
\)，故
\(
\mathcal E(a_1,\dots,a_T)=\mathcal E_{P_T,K(a)}
\)。
而 \(K\) 在 \(0\le k<P_T\) 上为双射，故像集为全体 \(\{\mathcal E_{P_T,k}\}\) 且对应唯一，从而单射成立。
\end{proof}

\begin{theorem}[交换外置的去序性与最小非交换补偿的初对象]\label{thm:conclusion-affine-ext-initial-object}
设 \(\Sigma\) 为满足 \(|\Sigma|\ge 2\) 的字母表，\(\Sigma^\ast\) 为自由幺半群。
交换幺半群同态对序信息的不可恢复性由定理 \ref{thm:pom-commutative-prime-register-order-nonrecoverable} 给出。
为刻画“在保留可交换伸缩子幺半群的同时，以最小非交换自由度补偿次序”的结构，定义范畴 \(\mathbf{AffExt}\) 如下：
对象为三元组 \((S,\iota,\tau)\)，其中 \(S\) 为幺半群，\(\iota:(\NN_{>0},\times)\hookrightarrow S\) 为单射幺半群同态且 \(\iota(\NN_{>0})\) 交换，\(\tau\in S\) 满足伸缩关系
\[
\iota(n)\tau=\tau^{n}\iota(n)\qquad(\forall\,n\in\NN_{>0}),
\]
态射为保持 \(\iota\) 与 \(\tau\) 的幺半群同态。
则 \(\bigl(\NN_{>0}\ltimes\NN,\iota_0,\tau_0\bigr)\) 为 \(\mathbf{AffExt}\) 的初对象，其中
\[
\iota_0(n)=(n,0),\qquad \tau_0=(1,1),
\qquad (N,k)(M,\ell)=(NM,k+N\ell).
\]
\leanverified{paper\_conclusion\_affine\_ext\_initial\_object}
\end{theorem}
\begin{proof}
对任意对象 \((S,\iota,\tau)\)，定义
\[
F:\NN_{>0}\ltimes\NN\to S,\qquad F(N,k):=\tau^{k}\iota(N).
\]
利用关系 \(\iota(N)\tau^\ell=\tau^{N\ell}\iota(N)\)（对 \(\ell\) 归纳得到），有
\[
F(N,k)F(M,\ell)
=\tau^k\iota(N)\tau^\ell\iota(M)
=\tau^{k+N\ell}\iota(NM)
=F(NM,k+N\ell),
\]
故 \(F\) 为幺半群同态，并满足 \(F\circ\iota_0=\iota\)、\(F(\tau_0)=\tau\)。
若 \(F'\) 亦保持 \(\iota_0,\tau_0\)，则 \(F'(N,0)=\iota(N)\)、\(F'(1,1)=\tau\)，从而 \(F'(N,k)=\tau^k\iota(N)\) 被唯一确定，故初对象成立。
\end{proof}
